\numberlesssection{РЕФЕРАТ}
Расчётно-пояснительная записка \pageref{LastPage} страница, \total{figures} рисунков, 
\total{tables} таблиц, \total{references} источник.

\noindent
МИКРОСЕРВИСНАЯ АРХИТЕКТУРА, ОТКАЗОУСТОЙЧИВОСТЬ, \newline ГРАФОВАЯ СУБД, РЕЛЯЦИОННАЯ СУБД, 
АЛГОРИТМЫ ОБХОДА ГРАФА, ПРОАКТИВНЫЙ АНАЛИЗ, NEO4J, POSTGRESQL, C\#, .NET 

Объектом исследования является топология ИТ-инфраструктуры распределенных 
информационных систем на основе микросервисной архитектуры.
Объектом разработки является база данных и прикладное программное обеспечение 
для анализа отказоустойчивости и моделирования каскадных сбоев в распределенных 
системах.

Цель работы --- проектирование и разработка базы данных, поддерживающей  
хранение высокосвязных данных и обеспечивающей высокую скорость выполнения 
алгоритмов обхода графа для поиска критических уязвимостей.

В аналитическом разделе проведен анализ предметной области управления надежностью 
систем и существующих коммерческих решений, показавший необходимость разработки 
проактивного инструмента для анализа топологии на этапе проектирования. 
Сформулированы требования к данным и ролевой модели пользователей. На основе 
анализа выбрана концепция полиглотного хранения данных, сочетающая реляционную 
и графовую модели.

В конструкторском разделе спроектирована структура хранения реляционной и графовой 
баз данных. На формальном языке описаны ограничения целостности и ролевая модель 
разграничения прав доступа на уровне приложения. Проведена проверка нахождения 
спроектированной реляционной базы данных в третьей нормальной форме.

В технологическом разделе обоснован выбор СУБД PostgreSQL и Neo4j, а также языковых 
средств разработки: платформы .NET 10, языка C\# и объектно-реляционного отображения 
Entity Framework Core. Представлены исходные коды описания схем данных, а также 
тексты спроектированных запросов с использованием библиотеки процедур APOC. 
Описана методика тестирования и приведены примеры работы интерфейса приложения.

В исследовательском разделе проведена практическая оценка производительности 
графовых запросов. Исследована зависимость времени выполнения анализа каскадных 
отказов от общего числа узлов системы, а также зависимость времени поиска 
циклических зависимостей от количества циклов в графе. Определено, что время 
выполнения сложных аналитических расчетов на топологиях объемом до десяти 
тысяч узлов не превышает нескольких десятков миллисекунд.